\chapter{Customer Partnership}

Customer Partnership represents an important and strategic pillar to accomplish the Bank's goals. The IT Department must establish measurable relationships between external customers and other stakeholders. Assuming a role that actively contributes to the improvement of the services provided and, in specific, monitors through KPIs the value perceived by customers.
The IT has to collaborate also internally within the organization in order to engage the appropriate department whenever it is needed.

This sector is extremely important for the success of the Bank's vision, because the international expansion through acquisitions and other partnerships introduces new business requirements due to different customer expectations, cultural diversity, and increased service complexity.
These challenges must be supported in a proactive way, by ensuring that IT services remain aligned with business objectives, deliver consistent value across different countries, and contribute to the Bank's growth; while complying with agreed Service Level Agreements between customers, stakeholders and user expectations.


These requisites make performance measurements crucial. By continuously monitoring and collecting customer-oriented metrics and analyzing services scores, IT Department is able to implement a continual service improvements that can strengthen customer trust and enrich the service quality while gaining long-term partneship. \\
This customer-centric approach supports organization`s objectives of delivering high-performance, high-quality and cost-effective IT services, aimed toward to accomplish the predefined goals previously mentioned in Cap.\ref{sec:goals}


\section{Time To Market}
\label{subs:timetomarket}

Time to Market, also known as ``Speed to Market'', is the organization's capability to deliver new or improved IT services in the shortest possible time, concerning from the concept to production lifecycle, so reflects the efficiency of the service design, development, testing, and deployment processes.

Reducing Time to Market allows the bank to respond quickly to changing business needs and market opportunities; in way to support digital banking services and international expansion. The objective is to accelerate service delivery without compromising quality, security or regulatory compliance.\\
Achieving this CSF requires effective collaboration between business and IT, supported by well-defined ITSM processes. \\
In the end, a shorter Time to Market contributes to increased customer satisfaction and competitive advantage.



\subsection{Average Lead Time for Change to Production}
\label{sub:AverageLeadTime}
Time to market is made operational through \textbf{Average Lead Time for Change to Production} KPI metric, which defines as the elapsed time between formal approval of a business requirement and its deployment into the production environment. \\
This metric is selected because it directly reflects the organization's ability to translate customer and business demands into operational banking capabilities which is essential for achieving the stated rapid digital service evolution and revenue growth objectives of the organisation.\\ 
In the context of the strategic goal of the organisation is to increase the income from internet banking, reduced Time to Market provides higher organizational agility and faster delivery of customer targeting features, thus improving market responsiveness and competitive positioning.

There are several ITIL aligned ITSM processes that primarily effects Time to Market:

\begin{itemize}
    \item \textbf{Change enablement} has a direct effect through approval workflows and Risk Management procedures, were excessive governance overhead can result in significant delays.
    \item \textbf{Release and Deployment Management} also effects the indicator  as deployment scheduling , release batching and automation maturity determine how quickly the validated changes end up deployed in the production.
    \item \textbf{Service Validation and Testing} influences Lead Time through test cycle duration and error handling loops, which can create bottlenecks if automation coverage is insufficient.
    \item \textbf{Service Design and Transition} activities also effect indirectly throught the time required for requirement analysis, architectural validation, and coordination of technical dependencies.
\end{itemize}

Overall, inefficiencies or lack of cooperation across these processes will result in the increase of Lead Time, where higher process integration, standardization and automation would reduce it, resulting with the improvement of Time To Market KPI.


% COMMETED SINCE PASTED FROM 3CSF
% % Average Lead Time of Change to Production measures the average amount of time taken for a new or modified IT service, feature, or change to progress from formal approval to successful deployment into the production environment. The metric evaluates the IT department's ability to rapidly deliver business value and respond to customer and market requirements. The reduction in lead time indicates increased organizational agility, improved collaboration between business and IT, and greater effectiveness of the service delivery process. \\

% % The Average Lead Time of Change to Production was selected because it directly measures the bank's ability to rapidly deliver new and modified digital services to customers. Considering the organization's strategic objectives of international expansion and increasing internet banking revenues, reducing the lead time of service deployment contributes directly to customer satisfaction, competitive advantage, and business agility. This metric is primarily influenced by the Service Design, Change Enablement, Release and Deployment Management, and Service Validation and Testing processes. \\

\begin{table}[h]
    \centering
    \renewcommand{\arraystretch}{1.6}
    \setlength{\arrayrulewidth}{0.6pt}

    \begin{tabularx}{\textwidth}{|X|X|}
    \hline
    \multicolumn{2}{|c|}{\cellcolor{mcblue}\textcolor{white}{\textbf{Average Lead Time for Change to Production}}} \\ \hline

    \rowcolor{mclightgray}
    \textbf{Description} & \textbf{Unit} \\ \hline
    Average time required for a new or modified IT service, feature or change to progress from a formal approval to successful deployment to the production. &
    Days. \\ \hline

    \rowcolor{mclightgray}
    \textbf{Formula} & \textbf{Frequency} \\ \hline
    \[
       \frac{\sum (\text{Deployment Date} - \text{Approval Date})}{\text{\# of Successfully Implemented Changes}}
    \]
     &
    Monthly, Quarterly. \\ \hline

    \end{tabularx}

    \caption{Time To Market - Avg Lead Time for Change to Production}
    \label{tab:measurement_AverageLeadTimeforChangetoProduction}
\end{table}


%   Savas

% Time to market, also known as ``Speed to Market'' is made operational through \textbf{Average Lead Time for Change to Production} KPI metric, which defines as the elapsed time between formal approval of a business requirement and its deployment into the production environment. \\
% This metric is selected because it directly reflects the organization's ability to translate customer and business demands into operational banking capabilities which is essential for achieving the stated rapid digital service evolution and revenue growth objectives of the organisation. In the context of the strategic goal of the organisation is to increase the income from internet banking, reduced Time to Market provides higher organizational agility and faster delivery of customer targeting features, thus improving market responsiveness and competitive positioning.

% There are several ITIL aligned ITSM processes that primarily effects Time to Market:

% \begin{itemize}
%     \item \textbf{Change enablement} has a direct effect through approval workflows and Risk Management procedures, were excessive governance overhead can result in significant delays.
%     \item \textbf{Release and Deployment Management} also effects the indicator  as deployment scheduling , release batching and automation maturity determine how quickly the validated changes end up deployed in the production.
%     \item \textbf{Service Validation and Testing} influences Lead Time through test cycle duration and error handling loops, which can create bottlenecks if automation coverage is insufficient.
%     \item \textbf{Service Design and Transition} activities also effect indirectly throught the time required for requirement analysis, architectural validation, and coordination of technical dependencies.
% \end{itemize}

% Overall, inefficiencies or lack of cooperation across these processes will result in the increase of Lead Time, where higher process integration, standardization and automation would reduce it, resulting with the improvement of Time To Market KPI.

% SHIFTED ON CHAP 4


%   DONE but check

\section{Performance To Contract}
\label{subs:performancetocontract}
% PerformanceToContract --> Alberto

Performance to Contract focuses on the IT Department's ability to respect contractual obligations and consistently deliver services according to the agreements made with customers and explained in natural language in the Service Level Agreements (\gls{sla}s). \\

This subject requires IT services to meet the agreed levels by different measurable values, such as availability, reliability, responsiveness, and quality; while ensuring and supporting the Bank's operational and predefined business objectives.


% TODO - make a list of definitions, to take more space?? 
% \begin{itemize}
%     \item \textbf{availability}: 
%     \item \textbf{reliability}:
%     \item \textbf{responsiveness}:
%     \item \textbf{quality}: 
    
% \end{itemize}

Well-defined Key Performance Indicators ensure evidence of whether service targets are being achieved and enable the IT department to identify eventual deviations or deficiencies. This gives a clear input to potential corrective actions or improvements to service performance over time. Moreover, the monitoring of performance against contractual commitments also increases transparency and accountability between IT and different business units or external third-party service providers.


This process is fundamental and very important for the Bank, because many banking services, such as Online Banking, ATMs, global payments, or other third-party services, depend strictly on service levels and regulatory compliance. Also, through a cybersecurity point of view, the respect of performance is fundamental to prevent various cyber-attacks and potential disservices.\\
In addition, being a global Bank, the services must be ensured not only during the local time of the main Headquarters, but, the support must be available respecting different cultural scenarios, thus supporting suppliers or international partners.

As a result of continuously measuring contractual performance and acting on the collected data, IT Department can improve service reliability, reduce the risk of SLA breaches and improve the overall customer relationships and partnerships with various stakeholders.



\subsection{SLA Compliance Rate}
\label{sub:slaComplianceRate}

SLA Compliance Rate can provide a measure whether contractual service levels are respected, thus to accomplish the agreements made with customers or stakeholders and therefore maintaint the quality of the service provided.


For this ~\gls{kpi} the process affected are:
\begin{itemize}
    \item \textbf{Service Level Management}: to ensure an opportune level of service provided that should ensure what agreed in SLA.
    \item \textbf{Incident Management}: to ensure that eventual incident are resolved and the serivce restored to its normal operativity and agreed quality.
    \item \textbf{Request Fulfillment}: to gather and fulfill customer and user requests and making it sure to accomplish them
    \item \textbf{Availabiliy Management}: to ensure the service is available for the user and the downtime dued to incidents or maintenance do not take more than necessary resources (time and costs).\\ 
    The availability of the service can be expressed via the forumla: 
    \[
        \text{Availabiliy} = \frac{\text{AST - downtime}}{AST} \times 100
    \]
    Where ``AST'' indicate the Agreed Service Time, which measure the expectation of the time that service is available; downtime should be included within this value.
    \item \textbf{Capacity Management}: to monitor and ensure that service achieve agreed and exprected performance, satisfying current and future demand in a cost-effective way.
    \item \textbf{Supplier Management}: To ensure that the organization`s suppliers and their performances are managed appropriately to support the seamless provision of quality products and services.
\end{itemize}


\begin{table}[h]
    \centering
    \renewcommand{\arraystretch}{1.6}
    \setlength{\arrayrulewidth}{0.6pt}

    \begin{tabularx}{\textwidth}{|X|X|}
    \hline
    \multicolumn{2}{|c|}{\cellcolor{mcblue}\textcolor{white}{\textbf{SLA Compliance Rate}}} \\ \hline

    \rowcolor{mclightgray}
    \textbf{Description} & \textbf{Unit} \\ \hline
    Percentage of request solved respecting SLAs specification, where is desired a high percentage. &
    Percentage. \\ \hline

    \rowcolor{mclightgray}
    \textbf{Formula} & \textbf{Frequency} \\ \hline
    \[
        \frac{\text{Requests solved within SLA}}{\text{Total requests}} \times 100
    \]
     &
    Monthly. \\ \hline

    \end{tabularx}

    \caption{Performance To Contract - SLA  Compliance Rate}
    \label{tab:measurement_SLAComplianceRate}
\end{table}


\subsection{Mean time for restoring service}
\label{sub:meantTimeRestore}
Another Key Performance Indicator which can help the IT Department toward the monitor service performance within the expectation of the stakeholders is the measurement of the restoring service time.\\
This indicator measures how quickly the IT organization restores service after an interruption, closely related to the previous metric~\ref{sub:slaComplianceRate} but can be adopted by IT Department as an input for continuous improvement toward to optimize its response for incidents.

Primarily involves the \textbf{``Incident Management''} process but can be extended to others related process such as:
\begin{itemize}
    \item \textbf{Problem Management} which cover the role to resolve the problems that caused the incident
    \item \textbf{Service Desks} which represent a single point of contact for users and other stakeholders, incident reports should pass by this department, which can be involved within the incident managament and classification then consume time and enter within the metric.
    \item \textbf{Service Configuration Management} (SACM) since accurate configuration information helps support teams to quickly identify affected Configuration Items (\gls{ci}s) and their dependencies, reducing diagnosis time. 
\end{itemize}

\begin{table}[H]
    \centering
    \renewcommand{\arraystretch}{1.6}
    \setlength{\arrayrulewidth}{0.6pt}

    \begin{tabularx}{\textwidth}{|X|X|}
    \hline
    \multicolumn{2}{|c|}{\cellcolor{mcblue}\textcolor{white}{\textbf{Mean time for restoring service}}} \\ \hline

    \rowcolor{mclightgray}
    \textbf{Description} & \textbf{Unit} \\ \hline
    Provide the average time required to restore an IT service after an incident has been reported. Lower values indicate better operational perforance and higher compliace with contractual SLAs.
     &
    Minutes. \\ \hline

    \rowcolor{mclightgray}
    \textbf{Formula} & \textbf{Frequency} \\ \hline
    \[
     \frac{\sum{\text{Restoration time}}}{\text{Number of resolved incidents}}
    \]
    
     &
    Weekly (operational team) \& Montly (as management reporting). \\ \hline

    \end{tabularx}

    \caption{Backlog - Priority Distribution}
    \label{tab:measurement_meanTimeRestoringService}
\end{table}



% ---------------------------------------------------------

%%%%%%%%%%%%%%%%%%%%%%%%%%%%%%%%%%%%%%%%%%%%%%%%%%%%%%%%%%%%%%%%%%%%%%%%%%%%%%%%%%%%%%%%%%%%%%%%%%%%%%%%%%%
%%%%%%%%%%%%%%%%%%%%%%%%%%%%%%%%%%%%%%%%%%%%%%%%%%%%%%%%%%%%%%%%%%%%%%%%%%%%%%%%%%%%%%%%%%%%%%%%%%%%%%%%%%%

